% Part: incompleteness
% Chapter: theories-computability
% Section: consis-with-q

\documentclass[../../../include/open-logic-section]{subfiles}

\begin{document}

\olfileid{inc}{tcp}{con}

\olsection{النظريات المتسقة مع $\Th{Q}$ غير قابلة للقرار}

لا تقف النتيجة الآتية عند عدم قابلية $\Th{Q}$ للقرار؛ بل تتناول
كل نظرية لا تناقض $\Th{Q}$، فتحكم بعدم قابليتها للقرار أيضًا.

\begin{thm}
إذا كانت $\Th{T}$ نظرية في لغة الحساب متسقة مع $\Th{Q}$، أي إن
$\Th{T}\cup\Th{Q}$ متسقة، فإن $\Th{T}$ غير قابلة للقرار.
\end{thm}

\begin{proof}
للنظرية $\Th{Q}$ بديهيات معدودة عددًا منتهيًا، هي $!Q_\OLMathNumeral{1}$، \dots،
$!Q_\OLMathNumeral{8}$. ويمكن جمعها في بديهية واحدة:
$!E=!Q_\OLMathNumeral{1}\land\dots\land !Q_\OLMathNumeral{8}$.

لنفرض $\Th{T}$ نظرية قابلة للقرار ومتسقة مع $\Th{Q}$، ولنضع
\[
\OLId{latin-looped}{C}=\Setabs{!A}{\Th{T}\Proves !E\lif !A}.
\]
نبيّن أن $\OLId{latin-looped}{C}$ تفصل حسابيًا بين $\Th{Q}$ و
$\Th{\bar Q}$، فنبلغ التناقض. أما إذا انتمت $!A$ إلى $\Th{Q}$،
فإن $!A$ تثبت من بديهيات $\Th{Q}$؛ فبمبرهنة الاستنباط يوجد
!!a{derivation} لـ$!E\lif !A$ في منطق الرتبة الأولى.
إذن $!A$ عنصر في $\OLId{latin-looped}{C}$.

وأما إذا انتمت $!A$ إلى $\Th{\bar Q}$، فثمة برهان على
$!E\lif\lnot !A$ في منطق الرتبة الأولى. ولو أثبتت $\Th{T}$
كذلك $!E\lif !A$، لأثبتت $\Th{T}$ نفي البديهية، أي $\lnot !E$؛
وحينئذ لا تكون $\Th{T}\cup\Th{Q}$ متسقة. ولكن
$\Th{T}\cup\Th{Q}$ متسقة بالفرض، فيمتنع أن تثبت $\Th{T}$
الجملة $!E\lif !A$، وتمتنع عضوية $!A$ في $\OLId{latin-looped}{C}$.

فكل $!A$ في $\Th{Q}$ تقع في $\OLId{latin-looped}{C}$،
وكل $!A$ في $\Th{\bar Q}$ تقع في $\Complement{\OLId{latin-looped}{C}}$.
وهكذا تكون $\OLId{latin-looped}{C}$ فاصلًا قابلًا للحساب، وهو التناقض المطلوب.
\end{proof}

ومن سعة هذه المبرهنة أنها تعطينا، مثلًا، النتيجة الآتية:

\begin{cor}
  في لغة الحساب لا يقبل منطق الرتبة الأولى القرار؛ والمقصود مجموعته
  $\Setabs{!A}{\text{الصيغة $!A$ قابلة للإثبات في منطق الرتبة الأولى}}$.
\end{cor}

\begin{proof}
ذلك أن منطق الرتبة الأولى مجموعة ما يلزم عن $\emptyset$،
وهذه المجموعة متسقة مع $\Th{Q}$.
\end{proof}

\end{document}
